\subsection{Computational certificate for $\mathrm{VDepth}(4,9)=7$}

\paragraph{Statement.}
For V-type 4-selection networks on 9 wires,
\[
  \mathrm{VDepth}(4,9)=7 .
\]
Here depth means the number of comparator layers. Each layer is a matching of
disjoint comparators.

\paragraph{Model and verification convention.}
A comparator network on \(n\) wires is a sequence of layers, where each layer
contains pairwise-disjoint comparators. A V-type \(t\)-selection network is
required to place the \(t\) selected outputs in the prescribed V-prefix
positions for every input. By the 0--1 principle, the network conditions are
checked on binary inputs. For the upper bound below, the stored SAT witness was
checked by an independent Python verifier on all \(2^9\) binary inputs.

\paragraph{Upper bound.}
The certificate includes an explicit depth-7 witness for the V-type
4-selection problem on 9 wires. The witness is stored at
\texttt{experiments/baseline\_n9/decisions/V\_t4\_n9/depth\_07/network.json}.
The verifier artifact
\texttt{experiments/baseline\_n9/decisions/V\_t4\_n9/depth\_07/verifier.json}
records \texttt{ok=true}, \texttt, \texttt{t=4}, and \texttt{n=9},
with the message that all binary inputs were verified. Thus the upper bound
\(\mathrm{VDepth}(4,9) \le 7\) is certified by an explicit network and
exhaustive binary simulation.

\paragraph{Lower bound strategy.}
The lower bound is organized as a dependency graph whose machine-readable form
is \texttt{experiments/v49\_d6\_final/proof\_dependency\_graph.json}. The graph
separates mathematical lemmas, artifact audits, solver UNSAT artifacts, and
derived conclusions.

First, the raw solver artifact
\texttt{experiments/baseline\_n9/decisions/V\_t3\_n9/depth\_05/solver.json}
records \texttt{UNSATISFIABLE} for \(V(3,9)\) at depth 5; this is a
\texttt{SOLVER\_UNSAT} premise. Second, the prefix observation that every
V-type 4-selection network is also a V-type 3-selection network is recorded as
a \texttt{MATHEMATICAL\_LEMMA}. Together with the empty-layer elimination
lemma, this implies that a depth-6 \(V(4,9)\) candidate with an empty layer
would yield a depth-at-most-5 \(V(3,9)\) candidate. Hence the depth-6 search may
be restricted to all layers nonempty; this is a \texttt{DERIVED\_CONCLUSION}.

The remaining depth-6 search is then reduced by the closure audit and
terminal-profile reduction, recorded as \texttt{ARTIFACT\_AUDIT} nodes. These
audits leave three residual parent families. The L5 tail census exhaustively
partitions these residual parents into 75 tail categories. Finally, for each of
the 75 leaves, the certificate records a raw clingo JSON artifact with
\texttt{Result=UNSATISFIABLE}; these leaf closures are \texttt{SOLVER\_UNSAT}
evidence for the corresponding constrained ASP instances. Combining the
audited coverage with the 75 raw UNSAT leaf artifacts gives the derived
depth-6 conclusion that no V-type 4-selection network on 9 wires exists at
depth 6. Figure~X schematically summarizes this partition.

\paragraph{Tail partition summary.}
The following table is read from
\texttt{experiments/v49\_d6\_final/tail\_leaf\_table.csv} and
\texttt{experiments/v49\_d6\_final/tail\_family\_summary.csv}.

\begin{center}
\begin{tabular}{lrrrrr}
\hline
Residual parent & Categories & L5 matchings & Screening & Overnight & UNSAT leaves \\
\hline
\texttt{A5\_residual\_L5eq3} & 9 & 1060 & 2 & 7 & 9 \\
\texttt{B6\_residual\_L5eq3} & 57 & 974 & 43 & 14 & 57 \\
\texttt{B6\_residual\_L5eq4} & 9 & 672 & 4 & 5 & 9 \\
\hline
Total & 75 & 2706 & 49 & 26 & 75 \\
\hline
\end{tabular}
\end{center}

The screening/overnight distinction records when the raw UNSAT artifact was
obtained; it is not used as a mathematical explanation of UNSAT.

\paragraph{Reproducibility and certificate.}
The final certificate is
\texttt{experiments/v49\_d6\_final/final\_certificate.json}. The evidence
manifest is \texttt{experiments/v49\_d6\_final/evidence\_manifest.json}, and
the evidence index is
\texttt{experiments/v49\_d6\_final/evidence\_index.md}. The dependency graph is
\texttt{experiments/v49\_d6\_final/proof\_dependency\_graph.json}. The audit
and bundling script is \texttt{scripts/audit\_v49\_d6\_final\_bundle.py}; the
tail-leaf analysis script is \texttt{src/analyze\_v49\_d6\_tail\_leaves.py}.
Raw solver outputs are included through the manifest, and SHA-256 hashes for
the source snapshots and evidence files are recorded there.

\paragraph{Limitations of the certificate.}
The solver UNSAT outputs are not converted here into human-readable proofs, and
no UNSAT cores are presented. The current certificate also does not include an
independent second encoding of the depth-6 lower-bound search. Timing data are
operational observations only. Changing solver version or hardware may affect
runtime, but not the recorded branch definitions or the manifest checksums. A
future clean replay and an independent CNF/SAT encoding would strengthen the
certificate.

\paragraph{Suggested citation in the paper.}
By the computational certificate described in Appendix~X, no depth-six V-type
4-selection network on nine wires exists, while a depth-seven network is
verified by exhaustive binary simulation. Hence \(\mathrm{VDepth}(4,9)=7\).

The case \(V(4,9)\) is settled by a computational certificate: the depth-six
search is reduced to 75 certified UNSAT tail leaves, and a depth-seven witness
is independently verified on all binary inputs.

Appendix~X records the computational certificate for \(V(4,9)\), separating the
mathematical reductions, artifact audits, raw solver UNSAT leaves, and the
verified depth-seven witness.

